\documentclass[journal]{IEEEtran}

\usepackage[UTF8]{ctex}
\usepackage{cite}
\usepackage{amsmath,amssymb,amsfonts}
\usepackage{graphicx}
\usepackage{booktabs}
\usepackage{url}
\usepackage{xcolor}
\usepackage{microtype}
\usepackage{tikz}
\usetikzlibrary{arrows.meta,positioning,fit}
\renewcommand{\ttdefault}{cmtt}
\providecommand{\href}[2]{\url{#1}}
\def\UrlBreaks{\do\/\do-\do\_\do\.\do\:\do\?\do\=\do\&}
\Urlmuskip=0mu plus 1mu\relax
\makeatletter
\setlength{\@fptop}{0pt}
\setlength{\@dblfptop}{0pt}
\makeatother
\setlength{\abovedisplayskip}{5pt plus 2pt minus 2pt}
\setlength{\belowdisplayskip}{5pt plus 2pt minus 2pt}
\setlength{\abovedisplayshortskip}{3pt plus 2pt minus 2pt}
\setlength{\belowdisplayshortskip}{5pt plus 2pt minus 2pt}

\graphicspath{{../}}
\newcommand{\best}[1]{\textbf{#1}}

\title{面向身份保持人脸恢复的\\双先验冲突感知适配器研究}

\author{张智洋~(Zhiy Zhang)\\
武汉大学国家网络安全学院\\
空天信息安全与可信计算教育部重点实验室\\
Email: \texttt{zhiyzhang@whu.edu.cn}}

\markboth{IEEE Transactions on Image Processing，~Vol.~XX, No.~XX, 2026}%
{Zhang: ConsistentFace}

\begin{document}
\bstctlcite{IEEEexample:BSTcontrol}
\maketitle

\begin{abstract}
真实场景中的人脸图像在拍摄、存储和传播过程中常受到模糊、噪声、下采样和压缩等复杂退化，导致面部结构与细节受损。盲人脸恢复旨在从这类退化观测中重建高质量人脸。现有方法虽然能够改善视觉质量，但严重退化会减少可辨认的身份信息，使恢复结果仍可能偏离目标身份。参考图像和文本提示可分别补充身份与语义信息；这些外部条件只有与目标观测相符时才能提供有效指导，而过时或失配的参考、错误文本以及局部外观冲突均可能干扰恢复。针对这一问题，本文提出 ConsistentFace，一种用于身份保持人脸恢复的可信度分配适配器。ConsistentFace 以可靠退化观测为锚点，核验观测、参考集和文本之间的一致性，并结合退化强度估计参考与文本的相对可信度。模型先在整图层面分配参考、文本和回退条件，以处理外部条件整体失效；再围绕头发、眼镜/眼睛、胡须/嘴巴、年龄/皮肤和配饰/遮挡细化局部条件，以限制局部冲突的影响范围。训练阶段，冲突类型仅用于监督条件分配；推理阶段无须提供冲突标签。大量实验和人工评审表明，ConsistentFace 在保持恢复质量的同时，提高了条件冲突下的身份保真度与外观一致性，并对参考失效、文本错误和局部冲突具有良好的鲁棒性。
\end{abstract}

\begin{IEEEkeywords}
盲人脸恢复，图像增强，扩散模型，多模态条件控制，身份保真度，可信度分配。
\end{IEEEkeywords}

\noindent\textbf{EDICS:} 图像与视频的恢复与增强；图像与视频处理中的机器学习；图像与视频的生物特征分析。

\section{绪论}
\label{sec:main_intro}

真实场景中的人脸图像在拍摄、存储和传输过程中常受到模糊、噪声、下采样和压缩等复杂退化 \cite{wang2021gfpgan,wang2023dr2}。盲人脸恢复旨在从受到未知退化影响的低质量观测中恢复高质量人脸 \cite{wang2021gfpgan,wang2022restoreformer}。现有方法能够改善感知质量和面部细节，但严重退化会减少可用于辨认目标身份的信息，使恢复结果仍可能偏离真实身份 \cite{suin2024diffuse,liu2025faceme,niu2026iconface}。

为补充退化观测中减少的信息，现有研究主要采用生成先验和外部条件两类途径。预训练人脸生成器、离散码本和扩散模型为结构与细节恢复提供生成先验 \cite{wang2021gfpgan,zhou2022codeformer,lin2024diffbir}；同身份参考图像补充身份信息 \cite{li2018gfrnet,li2020asffnet,liu2025faceme}，文本提示则提供语义描述和可控恢复信息 \cite{xie2023tfrgan,li2025mcs,zhu2026a2bfr}。ControlNet、IP-Adapter 和 PAIR Diffusion 进一步通过空间控制、解耦交叉注意力和多模态引导，将图像与文本条件接入扩散模型 \cite{zhang2023controlnet,ye2023ipadapter,goel2024pairdiffusion}。上述方法解决了信息补充和条件接入问题，但当参考图像与文本提示同时参与恢复时，模型仍缺少判断二者是否适用于当前目标的机制。

这一缺口会使失效的外部信息进入恢复网络。参考图像可能因姿态、表情、光照、年龄变化或身份失配而不适用于当前人脸 \cite{li2018gfrnet,li2020asffnet,niu2026iconface}，自动描述模型和视觉语言模型也可能生成与图像内容不一致的文本 \cite{rohrbach2018object,li2023pope}；两类条件还可能仅在眼镜、胡须或遮挡等局部外观上发生冲突。退化观测则直接来自目标人脸，始终保留与目标的对应关系。模糊、噪声和压缩会减少其中可辨认的信息，却不改变它作为目标观测的可靠性 \cite{liu2025faceme,niu2026iconface}。由此，双先验恢复需要依次回答三个问题：参考与文本哪一方更符合目标观测，冲突影响整张人脸还是局部外观，以及两类外部条件均不可信时采用何种恢复路径。

为解决上述问题，本文提出 ConsistentFace，在扩散恢复主干的条件编码与去噪网络之间加入可信度分配适配器。给定退化观测 $y$、参考集合 $\mathcal{R}$ 和文本提示 $T$，该适配器以 $y$ 为可靠锚点，核验观测、参考集和文本之间的一致性，并结合退化强度估计参考与文本的相对可信度。估计结果分别驱动全局分配和局部分配：前者处理整图层面的参考或文本失效，后者围绕头发、眼镜/眼睛、胡须/嘴巴、年龄/皮肤和配饰/遮挡调整局部条件；当两类外部条件均缺少可信证据时，\mbox{模型使用回退条件}。整个过程不改变扩散主干结构。

图~\ref{fig:conflict_motivation} 展示了文本错误、错误或过时参考、局部外观冲突以及两类外部条件同时失效四种情形。直接拼接、平均或固定加权无法区分条件是否适用于当前目标，失效信息因而可能进入去噪网络并影响恢复结果。ConsistentFace 根据冲突来源和影响范围调整参考、文本与回退条件的贡献。

\begin{figure*}[!t]
\centering
\includegraphics[width=\textwidth]{figures/fig01_conflict_motivation.png}
\caption{双先验冲突下的身份保持人脸恢复问题。简单融合可能错误采纳失效参考或错误文本，并使局部冲突扩散到整体恢复结果；ConsistentFace 根据参考--文本匹配、观测核验和参考集稳定性等证据，先判断整图层面的条件可信度，再对局部外观属性进行细化分配。}
\label{fig:conflict_motivation}
\end{figure*}

实现上述条件分配需要解决三个相互关联的挑战。1）\textbf{冲突来源判别。} 较低的参考--文本相似度既可能由文本错误引起，也可能来自过时参考或身份失配；只有结合可靠观测和参考集内部关系，模型才能确定应降低哪一类条件的权重。2）\textbf{冲突范围判别。} 条件失效可能影响整张人脸，也可能仅涉及眼镜、胡须或遮挡；模型需要把局部冲突限制在对应外观属性内，避免形成整体条件抑制。3）\textbf{训练与推理一致性。} 冲突类型可用于构造训练监督，但推理时通常不可用；模型必须仅依据 $(y,\mathcal{R},T)$ 完成可信度判断和条件分配。

ConsistentFace 对上述挑战分别采用多源证据、两级分配和无标签推理。多源证据联合描述参考--文本匹配、外部条件与观测的一致性、参考集内部稳定性和退化强度，用于确定冲突来源；全局分配器与局部分配器分别处理整体失效和局部冲突，并在证据不足时启用回退条件；冲突类型只在训练阶段构造分配监督，推理阶段直接由输入证据产生条件权重。

\textbf{本文的主要贡献如下：}

\begin{itemize}
    \item \textbf{提出冲突感知双先验恢复设定：} 将可靠退化观测、参考图像集合和文本提示纳入统一恢复框架，以观测为锚点处理参考失效、文本错误、双条件失效和局部外观冲突。
    
    \item \textbf{设计两级可信度分配机制：} 依据多源一致性与退化证据估计参考和文本的相对可信度，在全局层面处理整体条件失效，在属性层面处理局部冲突，并在外部条件证据不足时使用回退条件。

    \item \textbf{构建条件分配的训练与评测方案：} 利用冲突类型构造排序与回退监督，并采用去相关约束降低提示模板和参考来源造成的数据捷径；推理阶段无须冲突标签。同一主干比较、反事实实验、多评估器验证和人工评审共同验证了方法的有效性。

\end{itemize}


\section{相关工作}
\label{sec:main_related}
\subsection{盲人脸恢复与生成先验}
盲人脸恢复通常利用生成先验补偿退化造成的结构与细节损失 \cite{wang2021gfpgan,zhou2022codeformer,lin2024diffbir}。GFPGAN 引入预训练人脸生成器先验 \cite{wang2021gfpgan}；RestoreFormer、VQFR 和 CodeFormer 分别利用高质量键值字典、向量量化字典和离散码本完成恢复 \cite{wang2022restoreformer,gu2022vqfr,zhou2022codeformer}。

近期方法进一步引入扩散先验。DR2、DifFace 和 DiffBIR 分别通过鲁棒退化去除、后验迭代恢复和潜空间生成先验改善恢复结果 \cite{wang2023dr2,yue2022difface,lin2024diffbir}。相关研究还将扩散先验用于真实图像超分辨率，并探索区域自适应、幻觉校正、视觉风格、3D 先验和参考梯度引导 \cite{wang2024stablesr,suin2024diffuse,gao2024diffmac,zhu2024visualstyle,lu2024threepriors,min2024pdgrad}。

生成先验主要从退化观测本身推断缺失内容。观测中的身份线索随退化增强而减少时，生成结果仍可能偏离目标身份 \cite{suin2024diffuse,liu2025faceme,niu2026iconface}。为补充这部分信息，后续研究开始引入外部身份或语义条件 \cite{li2018gfrnet,liu2025faceme,xie2023tfrgan,an2026tbfr}。

\subsection{参考驱动与个性化人脸恢复}
参考驱动方法通过额外的同身份图像补充身份信息。GFRNet 对引导图像进行变形，使其姿态和表情与退化观测对齐；ASFFNet 进一步结合空间对齐、光照变换和多参考特征自适应融合 \cite{li2018gfrnet,li2020asffnet}。DMDNet 则以通用字典和身份特定字典分别存储一般人脸先验与身份相关特征 \cite{li2022dmdnet}。这些方法表明，参考信息能否有效参与恢复取决于参考与观测之间的对齐和选择 \cite{li2018gfrnet,li2020asffnet,li2022dmdnet}。

个性化生成方法提供了另一种身份条件建模方式。MyStyle 学习面向特定身份的个性化生成先验 \cite{nitzan2022mystyle}；Textual Inversion、DreamBooth 和 Custom Diffusion 通过学习文本嵌入或调整模型参数表示特定主体 \cite{gal2023textualinversion,ruiz2023dreambooth,kumari2023customdiffusion}。IP-Adapter 提供通用图像提示接口，PhotoMaker、InstantID 和 PuLID 则进一步面向人像身份定制 \cite{ye2023ipadapter,li2024photomaker,wang2024instantid,guo2024pulid}。

FaceMe 将上述个性化条件用于盲人脸恢复：该方法使用身份编码器提取参考图像中的身份特征，并通过两阶段训练平衡低质量输入与参考图像之间的依赖，更换身份时无须重新微调 \cite{liu2025faceme}。这类方法通常以参考图像有效且与目标身份匹配为前提。ConsistentFace 在此基础上研究参考失效，以及参考图像与文本提示不一致时的条件选择。

\subsection{语言引导恢复与可变属性控制}
自然语言为恢复模型提供了面向用户的语义接口 \cite{conde2024instructir,qi2024spire}。TFRGAN 将文本与图像信息融合到 StyleGAN 潜空间，并利用文本特征调制解码过程 \cite{xie2023tfrgan}；TBFR 构建了自然语言引导的盲人脸恢复方法与数据集 \cite{an2026tbfr}。MCS 利用不同文本提示控制多样化恢复结果 \cite{li2025mcs}，A$^2$BFR 则通过图文交叉模态注意力和属性感知学习实现细粒度提示控制 \cite{zhu2026a2bfr}。这些研究从整体语义逐步扩展到局部属性，为文本参与人脸恢复提供了基础。

现有文本引导方法通常默认提示能够描述预期内容，但图像描述模型和大型视觉语言模型可能生成输入图像中不存在的对象 \cite{rohrbach2018object,li2023pope}，历史标注和用户描述也可能存在偏差。因此，文本条件参与恢复前还需要核验其与目标观测的一致性，并区分整体语义错误与局部属性冲突。

\subsection{多模态冲突与可信度分配}
图像与文本条件可以通过多种接口共同控制扩散模型。ControlNet 通过附加分支注入边缘、深度或姿态等空间条件 \cite{zhang2023controlnet}；IP-Adapter 在交叉注意力中分离图像提示和文本提示 \cite{ye2023ipadapter}；PAIR Diffusion 利用多模态无分类器引导组合参考图像与文本条件，支持对象级编辑 \cite{goel2024pairdiffusion}。这些方法主要回答不同条件如何接入和组合，尚未处理参考图像与文本提示冲突时的条件取舍。CLIP 表征可以衡量图文语义一致性 \cite{radford2021clip}，但单一相似度不足以区分文本错误、参考错误和双条件同时失效。

近期研究进一步协调身份、结构与文本条件。CoFaDiff 研究个性化生成中身份一致性与文本可编辑性的平衡 \cite{ming2026cofadiff}；IConFace 以参考图像提供全局身份信息，并利用退化图像保持空间结构，从而统一有参考与无参考人脸恢复 \cite{niu2026iconface}。与这些工作不同，本文将问题落在条件冲突上，依据可靠退化观测分配参考身份条件与文本条件的相对可信度。

\section{方法}
\label{sec:main_method}

\subsection{方法概述}
ConsistentFace 的总体框架由扩散恢复主干、观测控制分支和可信度分配适配器组成。恢复主干采用 Stable Diffusion XL（SDXL）\cite{podell2023sdxl}：其 VAE 将图像映射到低维潜空间，去噪网络在该空间内逐步恢复清晰潜变量 \cite{rombach2022ldm}。观测控制分支采用 ControlNet，将可靠退化输入 $y$ 编码为空间结构条件并注入去噪网络 \cite{zhang2023controlnet}。可信度分配适配器位于条件编码器与去噪网络之间，根据当前观测分配参考集合 $\mathcal{R}$、文本提示 $T$ 和回退条件的权重。

图~\ref{fig:overall_architecture} 按数据流给出五个区域。区域 1 接收可靠退化观测、参考集合和文本提示；区域 2 将三类输入编码为观测特征、参考条件 token 和文本语义 token；区域 3 计算参考--文本匹配、观测核验、参考集稳定性和退化强度；区域 4 根据这些证据生成全局条件与局部属性条件；区域 5 将分配后的条件 token 与 ControlNet 结构条件共同送入 SDXL 去噪网络，得到恢复结果。

区域 4 包含两个相互衔接的分配层级：全局分配器确定整张人脸上的参考--文本偏好和回退权重，局部属性分配器再围绕头发、眼镜/眼睛、胡须/嘴巴、年龄/皮肤和配饰/遮挡细化属性级权重。以下按照图中的数据流依次说明输入输出、条件编码、证据构造、两级分配以及训练与推理策略。

\begin{figure*}[!t]
\centering
\includegraphics[width=\textwidth]{figures/fig02_overall_architecture.png}
\caption{ConsistentFace 总体框架。模型先编码可靠退化观测、参考集合和文本提示，再构造多源一致性与退化证据。全局分配器和局部属性分配器据此生成条件 token，并与 ControlNet 提取的观测结构条件共同指导 SDXL 去噪。}
\label{fig:overall_architecture}
\end{figure*}

\subsection{问题定义}
设 $x$ 表示目标清晰人脸，$y=\mathcal{D}(x)$ 表示其可靠低质量观测，其中退化算子 $\mathcal{D}$ 可包含模糊、下采样、噪声和压缩等真实退化 \cite{wang2021gfpgan,wang2023dr2}。给定参考图像集合 $\mathcal{R}=\{R_i\}_{i=1}^{M}$ 和自然语言提示 $T$，人脸恢复模型生成结果
\begin{equation}
\label{eq:problem_restore}
\hat{x}=G_{\theta}(y,\mathcal{R},T).
\end{equation}
退化观测 $y$ 与目标人脸直接对应，是恢复过程始终保留的可靠输入；其结构与身份信息的可辨认程度由退化强度决定。参考集合 $\mathcal{R}$ 和文本提示 $T$ 提供额外的身份与外观信息，但可能因身份失配、时间变化、描述错误或局部属性差异而与 $y$ 不一致。因此，模型在使用两类外部条件前先估计其相对可信度。

双先验恢复由恢复分支和条件分配分支共同完成：前者生成 $\hat{x}$，后者估计参考、文本和回退条件的使用权重。记可信度分配器为
\begin{equation}
\label{eq:problem_allocator}
\{\alpha,c,\alpha_a,c_a\}_{a\in\mathcal{A}}
=A_{\psi}(y,\mathcal{R},T),
\end{equation}
其中，$\alpha$ 表示参考与文本的全局相对偏好；$c$ 表示外部条件整体使用权重，$1-c$ 表示全局回退权重。对于属性组 $a$，$\alpha_a$ 表示该组中参考与文本的局部相对偏好；$c_a$ 表示该组的外部条件使用权重，$1-c_a$ 表示局部回退权重。本文选取五类可解析外观属性构成 $\mathcal{A}$：头发、眼镜/眼睛、胡须/嘴巴、年龄/皮肤和配饰/遮挡，用于刻画局部冲突和细化条件分配；该集合不覆盖全部面部语义，也不被视为固定的最优属性划分。训练阶段的冲突类型仅用于构造分配监督，推理阶段模型直接由 $(y,\mathcal{R},T)$ 得到上述权重。

\subsection{条件编码}
图~\ref{fig:overall_architecture} 的区域 2 将三类输入转换为两组表示：观测、参考和文本的全局特征用于后续一致性计算，参考与文本的条件 token 则用于去噪网络。以下分别给出参考、文本和观测的编码过程。

对于参考集合 $\mathcal{R}=\{R_i\}_{i=1}^{M}$，模型先从每张图像提取身份特征和视觉特征，再由可学习映射器聚合为参考条件 token：
\begin{equation}
\label{eq:reference_tokens}
\begin{aligned}
e^{I}_{R_i}&=E_I(R_i),\\
Z_{R_i}&=E_V(R_i),\\
V&=\Phi_R\!\left(\{e^{I}_{R_i},Z_{R_i}\}_{i=1}^{M}\right),\\
V&=[v_1,\ldots,v_K].
\end{aligned}
\end{equation}
其中，$E_I$ 为身份编码器，本文采用 ArcFace \cite{deng2019arcface}；$E_V$ 复用 PhotoMaker 的 CLIP 视觉编码器 \cite{li2024photomaker}；$\Phi_R$ 为参考特征映射器。映射结果 $V=[v_1,\ldots,v_K]$ 是融合身份与视觉信息的参考条件序列。$K$ 由 $\Phi_R$ 的输出长度确定，本文设置为 8；文本分支输出相同数量和维度的 token，以便两类条件逐 token 组合。

文本分支沿用 SDXL 文本编码器 $E_T$，并增加映射器 $\Phi_T$ 以对齐参考条件空间。单条描述直接作为提示 $T$；多条描述按输入顺序以分隔符连接后仍记为 $T$，从而保持 SDXL 的标准文本输入形式。编码器输出序列特征 $H$，池化得到全局文本向量 $t$；映射器再将 $H$ 转换为 $K$ 个文本语义 token：
\begin{equation}
\label{eq:text_tokens}
\begin{aligned}
H&=E_T(T),\\
t&=\operatorname{Pool}(H),\\
U&=\Phi_T(H)=[u_1,\ldots,u_K].
\end{aligned}
\end{equation}
其中，$H$ 用于保留逐 token 语义，$t$ 用于计算全局一致性，$U=[u_1,\ldots,u_K]$ 用于条件融合。$\Phi_T$ 同时对齐 $U$ 与 $V$ 的维度和序列长度，使两类条件可以共享分配权重；SDXL 的原始文本编码能力保持不变。

观测分支对 $y$ 执行两种编码：ControlNet 提取供去噪网络使用的结构条件，观测编码器 $E_y$ 提取供可信度核验使用的特征 $f_y$。为与全局文本向量 $t$ 对应，模型同时从完整参考序列 $V$ 中计算全局参考摘要：
\begin{equation}
\label{eq:mean_tokens}
f_y=E_y(y),\qquad
\bar{v}=\frac{1}{K}\sum_{k=1}^{K}v_k.
\end{equation}
其中，$\bar{v}$ 只用于全局一致性计算。最终条件融合仍使用完整序列 $V$ 和 $U$，局部属性分配也从 $V$ 中单独提取属性相关特征，因此式~\eqref{eq:mean_tokens} 的均值操作不参与最终 token 融合。

\subsection{一致性与退化证据构造}
图~\ref{fig:overall_architecture} 的区域 3 将编码结果转换为分配证据。证据包括外部条件之间的一致性、外部条件与观测的一致性、参考集合内部一致性，以及反映观测可辨认程度的退化强度。前三类证据采用余弦相似度计算：
\begin{equation}
\label{eq:cosine_similarity}
\operatorname{sim}(a,b)=
\frac{a^{\top}b}{\|a\|_2\,\|b\|_2}.
\end{equation}
该相似度用于比较投影到同一空间的特征，取值越大表示二者越一致。

三类一致性关系具体写为四个全局分数：
\begin{subequations}
\label{eq:multi_evidence_scores}
\begin{align}
s_{RT}
&=\operatorname{sim}\!\left(P_R(\bar{v}),P_T(t)\right),\\
s_{yR}
&=\operatorname{sim}\!\left(P_y^{I}(f_y),P_R(\bar{v})\right),\\
s_{yT}
&=\operatorname{sim}\!\left(P_y^{A}(f_y),P_T^{A}(t)\right),\\
s_{RR}
&=\frac{2}{M(M-1)}
\sum_{1\leq i<j\leq M}
\operatorname{sim}\!\left(e^{I}_{R_i},e^{I}_{R_j}\right).
\end{align}
\end{subequations}
其中，$s_{RT}$ 衡量参考与文本的匹配程度，$s_{yR}$ 和 $s_{yT}$ 分别核验参考身份和文本外观是否符合可靠退化观测，$s_{RR}$ 衡量参考集内部身份是否稳定。$P_R$、$P_T$、$P_y^{I}$、$P_y^{A}$ 和 $P_T^{A}$ 为可学习投影头，用于建立可比较的特征空间。仅有一张参考图像时，$s_{RR}$ 取中性常数，避免不存在的集合关系影响分配。

观测中可辨认信息的多少会影响上述核验结果，因此模型进一步估计退化强度：
\begin{equation}
\label{eq:degradation_score}
q=Q(f_y).
\end{equation}
其中，$Q$ 是接在观测编码器 $E_y$ 之后的轻量可学习退化强度估计头，输出标量 $q$。$q$ 越大，表示可靠退化观测中可用于核验身份和属性的线索越弱；这一项刻画观测线索的可辨认程度，而不改变 $y$ 作为目标输入的可靠性。

基础证据由 $s_{RT}$、$s_{yR}$、$s_{yT}$、$s_{RR}$ 和 $q$ 构成。模型还计算 $s_{RT}-s_{yR}$ 和 $s_{RT}-s_{yT}$，用于比较跨先验不一致与观测对参考、文本的支持程度：在 $s_{RT}$ 接近时，前一差值较小表示观测更支持参考，后一差值较小表示观测更支持文本。最终，全局证据向量为
\begin{equation}
\label{eq:global_evidence_vector}
g=
\left[
s_{RT},\,
s_{yR},\,
s_{yT},\,
s_{RR},\,
q,\,
s_{RT}-s_{yR},\,
s_{RT}-s_{yT}
\right].
\end{equation}

\subsection{全局条件分配}
图~\ref{fig:reliability_allocation}A 展示全局分配过程。证据向量 $g$ 经过两个分配头后得到参考--文本偏好和外部条件使用权重：
\begin{subequations}
\label{eq:global_weights}
\begin{align}
\alpha&=\sigma\!\left(f_{\alpha}(g)\right),\\
c&=\sigma\!\left(f_{c}(g)\right).
\end{align}
\end{subequations}
其中，$f_{\alpha}$ 和 $f_c$ 为多层感知机，$\sigma(\cdot)$ 将输出限制在 $[0,1]$。$\alpha$ 控制参考条件 $V$ 与文本条件 $U$ 的相对贡献：$\alpha$ 越大，参考贡献越高；$c$ 控制两类外部条件与回退条件的相对贡献：$c$ 越小，回退贡献越高。

全局条件 token 因此写为
\begin{equation}
\label{eq:global_fusion}
F_g=
c\Big(\alpha V+(1-\alpha)U\Big)
+(1-c)B_g.
\end{equation}
其中，$B_g$ 是与 $V$ 和 $U$ 维度相同的可学习回退 token。式~\eqref{eq:global_fusion} 由 $\alpha$ 完成外部条件内部的选择，再由 $c$ 决定是否采用外部条件。

\subsection{局部属性条件分配}
全局条件 $F_g$ 作用于整张人脸。为把局部冲突限制在对应外观区域，图~\ref{fig:reliability_allocation}B 进一步对五类属性组 $\mathcal{A}$ 分配条件，即头发、眼镜/眼睛、胡须/嘴巴、年龄/皮肤和配饰/遮挡。对于属性 $a\in\mathcal{A}$，模型先从文本序列特征 $H$ 中提取局部语义表示：
\begin{equation}
\label{eq:slot_text_feature}
t_a=\Psi_a(H).
\end{equation}
其中，$\Psi_a$ 为属性 $a$ 的文本特征提取器，$t_a$ 表示对应的局部文本语义。

参考分支使用属性查询对完整序列 $V$ 执行缩放点积注意力，以提取同一属性的参考特征 \cite{vaswani2017attention}：
\begin{subequations}
\label{eq:slot_attention}
\begin{align}
\omega_{a,k}
&=
\frac{
\exp\!\left((W_Q e_a)^{\top}(W_K v_k)/\sqrt{d}\right)
}{
\sum_{\ell=1}^{K}
\exp\!\left((W_Q e_a)^{\top}(W_K v_{\ell})/\sqrt{d}\right)
},\\
v_a
&=
\sum_{k=1}^{K}\omega_{a,k}W_V v_k.
\end{align}
\end{subequations}
其中，$e_a$ 是属性查询向量，$W_Q$、$W_K$ 和 $W_V$ 分别为查询、键和值的线性投影矩阵，$v_a$ 表示与属性 $a$ 对应的参考视觉特征。

得到两类局部表示后，适配器计算属性级一致性：
\begin{subequations}
\label{eq:slot_scores}
\begin{align}
s_a
&=\operatorname{sim}\!\left(P_R^{a}(v_a),P_T^{a}(t_a)\right),\\
s_{ya}
&=\operatorname{sim}\!\left(P_y^{a}(f_y),P_T^{a}(t_a)\right).
\end{align}
\end{subequations}
其中，$s_a$ 比较参考与文本在属性 $a$ 上的描述，$s_{ya}$ 核验该属性文本是否符合可靠退化观测。二者与全局证据共同输入局部分配器。

局部分配器据此输出属性 $a$ 的参考--文本偏好和外部条件使用权重：
\begin{subequations}
\label{eq:slot_weights}
\begin{align}
g_a&=\left[g,\,s_a,\,s_{ya},\,\alpha,\,c\right],\\
\alpha_a&=\sigma\!\left(f_{\alpha}^{a}(g_a)\right),\\
c_a&=\sigma\!\left(f_{c}^{a}(g_a)\right).
\end{align}
\end{subequations}
其中，$g_a$ 连接全局证据、局部一致性分数和全局分配结果，使局部判断保持与整图条件一致。$\alpha_a$ 控制参考特征 $v_a$ 与文本特征 $t_a$ 的相对贡献，$c_a$ 控制外部条件与局部回退条件的相对贡献。

局部属性融合 token 定义为
\begin{equation}
\label{eq:slot_fusion}
F_a=
c_a\Big(\alpha_a v_a+(1-\alpha_a)t_a\Big)
+(1-c_a)B_a.
\end{equation}
其中，$B_a$ 为属性 $a$ 的回退 token。全局条件与五组局部条件最终拼接为
\begin{equation}
\label{eq:condition_tokens}
F=\operatorname{Concat}\!\left(F_g,\{F_a:a\in\mathcal{A}\}\right).
\end{equation}
其中，$F$ 为注入 SDXL 去噪网络的完整条件 token 集合。

图~\ref{fig:reliability_allocation} 右侧对应三种分配结果：文本可信度较低时提高参考贡献，参考可信度较低时提高文本贡献，两类外部条件证据均不足时提高回退贡献。

\begin{figure*}[!t]
\centering
\includegraphics[width=\textwidth]{figures/fig03_reliability_allocation.png}
\caption{两级可信度分配。全局分配器输出参考--文本偏好 $\alpha$ 和外部条件使用权重 $c$；局部属性分配器输出属性组内的参考--文本偏好 $\alpha_a$ 和外部条件使用权重 $c_a$。}
\label{fig:reliability_allocation}
\end{figure*}

\subsection{训练目标与推理策略}
训练同时优化恢复结果与条件分配。恢复分支采用扩散去噪损失；分配分支采用全局与局部排序监督、回退监督和去相关约束。按照扩散模型的噪声回归目标 \cite{ho2020ddpm,rombach2022ldm,podell2023sdxl}，设 $z_0$ 为清晰图像 $x$ 的 VAE 潜变量，$\tau$ 为扩散时间步，$\epsilon$ 为标准高斯噪声，则加噪潜变量为
\begin{equation}
\label{eq:noisy_latent}
z_{\tau}
=
\sqrt{\bar{\gamma}_{\tau}}\,z_0
+
\sqrt{1-\bar{\gamma}_{\tau}}\,\epsilon.
\end{equation}
其中，$\bar{\gamma}_{\tau}$ 表示扩散噪声调度在时间步 $\tau$ 的累计保留系数。

对应的去噪损失为
\begin{equation}
\label{eq:diffusion_loss}
\mathcal{L}_{\mathrm{diff}}
=
\mathbb{E}_{z_0,\epsilon,\tau}
\left[
\left\|
\epsilon-\epsilon_{\theta}(z_{\tau},\tau,C_y,F)
\right\|_2^2
\right].
\end{equation}
其中，$C_y$ 是由可靠退化输入 $y$ 提取的 ControlNet 结构条件，$F$ 是可信度分配后的条件 token。

训练阶段使用冲突类型 $r$ 构造权重排序对和回退目标。令全局权重集合为
\begin{equation}
\label{eq:global_weight_set}
\mathcal{W}_g=\{\alpha,\ 1-\alpha,\ c,\ 1-c\},
\end{equation}
其中四项分别对应参考偏好、文本偏好、外部条件使用权重和回退权重。每个冲突类型 $r$ 被映射为监督排序对 $\mathcal{P}_g(r)\subseteq\{(\alpha,1-\alpha),(1-\alpha,\alpha),(c,1-c),(1-c,c)\}$。文本错误样本要求参考偏好高于文本偏好，错误或过时参考样本要求文本偏好高于参考偏好，双条件失效样本要求回退权重高于外部条件使用权重。全局排序损失为
\begin{equation}
\label{eq:global_rank_loss}
\begin{aligned}
\mathcal{L}^{g}_{\mathrm{rank}}
=&
\frac{1}{|\mathcal{P}_g(r)|}
\sum_{(w^+,w^-)\in\mathcal{P}_g(r)}
\\[-1pt]
&\max\!\left(0,\ m_g-(w^+-w^-)\right).
\end{aligned}
\end{equation}
其中 $m_g$ 是全局排序间隔。该损失约束权重的相对方向和间隔，无须为各类冲突预设固定权重；当 $\mathcal{P}_g(r)$ 为空时，该项置为 $0$。

局部属性权重同样采用成对间隔排序损失。对于受冲突影响的属性 $a\in\mathcal{A}^{-}(r)$，标签 $r$ 指定参考、文本或回退权重的监督顺序。令局部权重集合为
\begin{equation}
\label{eq:slot_weight_set}
\mathcal{W}_a=\{\alpha_a,\ 1-\alpha_a,\ c_a,\ 1-c_a\}.
\end{equation}
局部排序对 $\mathcal{P}_a(r,a)$ 仅在 $a\in\mathcal{A}^{-}(r)$ 时启用；无冲突属性不施加显式排序约束，以保留仍可信的文本指导。局部排序损失为
\begin{equation}
\label{eq:slot_rank_loss}
\begin{aligned}
\mathcal{L}^{a}_{\mathrm{rank}}
=&
\frac{1}{|\mathcal{A}^{-}|}
\sum_{a\in\mathcal{A}^{-}}
\frac{1}{|\mathcal{P}_a(r,a)|}
\\[-1pt]
&\sum_{(w_a^+,w_a^-)\in\mathcal{P}_a(r,a)}
\max\!\left(0,\ m_a-(w_a^+-w_a^-)\right).
\end{aligned}
\end{equation}
其中 $m_a$ 是局部排序间隔。当 $\mathcal{A}^{-}$ 为空时，该项置为 $0$。

回退损失监督模型在全局和局部属性层面是否启用回退。设 $b_g(r)\in\{0,1\}$ 表示全局是否需要回退，$b_a(r,a)\in\{0,1\}$ 表示属性 $a$ 是否需要回退，则
\begin{equation}
\label{eq:fallback_loss}
\begin{aligned}
\mathcal{L}_{\mathrm{fb}}
=&
-b_g\log(1-c)-(1-b_g)\log c\\
&+
\frac{1}{|\mathcal{A}|}
\sum_{a\in\mathcal{A}}
\left[-b_a\log(1-c_a)-(1-b_a)\log c_a\right].
\end{aligned}
\end{equation}
实现中对 $\log(\cdot)$ 的输入加入 $10^{-6}$ 下界以避免数值不稳定。

排序与回退监督直接约束分配方向，去相关约束则减少模型利用非目标元信息形成数据捷径。令 $\tilde{g}$ 和 $\tilde{g}_a$ 表示批内标准化后的全局证据与局部属性证据，$\xi$ 为提示模板与参考来源标识的一热编码拼接，则
\begin{equation}
\label{eq:decorrelation_loss}
\mathcal{L}_{\mathrm{decorr}}
=
\left\|\operatorname{Cov}(\tilde{g},\xi)\right\|_F^2
+
\frac{1}{|\mathcal{A}|}
\sum_{a\in\mathcal{A}}
\left\|\operatorname{Cov}(\tilde{g}_a,\xi)\right\|_F^2.
\end{equation}
该约束用于减少条件分配对特定提示模板和参考来源的依赖。

完整训练目标为
\begin{equation}
\label{eq:training_objective}
\mathcal{L}
=
\mathcal{L}_{\mathrm{diff}}
+
\lambda_g\mathcal{L}^{g}_{\mathrm{rank}}
+
\lambda_a\mathcal{L}^{a}_{\mathrm{rank}}
+
\lambda_b\mathcal{L}_{\mathrm{fb}}
+
\lambda_d\mathcal{L}_{\mathrm{decorr}}.
\end{equation}
其中，$\lambda_g$、$\lambda_a$、$\lambda_b$ 和 $\lambda_d$ 分别控制全局排序损失、局部排序损失、回退损失和去相关约束。本文设置 $m_g=m_a=0.20$、$\lambda_g=\lambda_a=\lambda_b=0.05$ 和 $\lambda_d=0.02$。除移除相应损失项的消融变体外，其余 ConsistentFace 变体采用相同设置。式~\eqref{eq:training_objective} 用于阶段 1 和阶段 3；阶段 2 固定参考映射器与可信度分配适配器，仅以 $\mathcal{L}_{\mathrm{diff}}$ 更新 ControlNet。上述数值在验证集上确定，并在测试阶段保持不变。

图~\ref{fig:training_inference} 汇总训练与推理流程。训练时，冲突类型用于构造排序对和回退目标，提示模板与参考来源标识用于计算去相关约束；推理时不再输入这些标签或元信息，模型仅根据 $(y,\mathcal{R},T)$ 计算 $g$ 与 $g_a$，并生成最终条件 $F$。

\begin{figure*}[!t]
\centering
\includegraphics[width=\textwidth]{figures/fig04_training_inference.png}
\caption{条件分配监督与无标签推理。训练阶段利用冲突类型构造全局及局部排序监督和回退目标，并利用非目标元信息计算去相关约束；推理阶段仅由 $y$、$\mathcal{R}$ 和 $T$ 估计条件可信度。}
\label{fig:training_inference}
\end{figure*}

\section{实验}
\label{sec:main_experiments}

实验围绕四个问题展开：不同先验输入下的方法处于何种性能区间；在相同恢复主干、输入和训练预算下，可信度分配是否优于直接融合；多源证据、局部属性分配和回退条件分别发挥何种作用；所得结论能否在不同提示、退化类型、评估器和人工评审中保持一致。

\subsection{实验设置}
\subsubsection{训练数据与实现细节}
\label{subsec:training_implementation}
训练清单包含 70,032 条 FFHQ 目标记录 \cite{liu2025faceme}，每条记录均配有语义提示和由 FFHQRef 提取的同身份参考特征。所有图像缩放至 $512\times512$，每次随机采样不超过四张参考图像，提示模板为 ``A photo of face. \{semantic\}''。生成主干采用 RealVisXL V3.0，参考身份编码权重由 PhotoMaker v1 初始化，其余可训练模块按照各阶段配置进行优化。

模型采用三阶段优化。阶段 1 固定 VAE、两个文本编码器和主去噪 UNet，联合训练参考映射器 $\Phi_R$、ControlNet 和可信度分配适配器 $A_\psi$，建立条件表示与分配规则。阶段 2 从阶段 1 检查点初始化，固定 $\Phi_R$ 与 $A_\psi$，仅以扩散去噪损失更新 ControlNet，使恢复分支在既定条件分配下充分优化。阶段 3 保持 $\Phi_R$ 固定，以较低学习率联合微调 ControlNet 和 $A_\psi$，对齐恢复结果与条件分配。

阶段 2 和阶段 3 还设置概率为 $0.1$ 的无参考身份 token 分支。该分支保留文本条件和由退化观测驱动的 ControlNet，不执行参考身份 token 的特征替换，并将适配器辅助损失置零。因此，它训练的是退化观测--文本条件恢复，而非脱离观测的纯文本生成；阶段 1 不启用该分支。

表~\ref{tab:training_hyperparams} 汇总三个阶段的训练设置。T、R 和 B 分别表示文本错误、参考错误和双条件失效的采样概率，M 表示其余匹配样本，满足 $M=1-T-R-B$。局部属性冲突由属性组标注和监督掩码表示，参与局部排序与回退监督，不另设与 T/R/B 并列的全局采样概率。

\begin{table*}[!t]
\centering
\caption{三阶段实际训练设置。全部阶段均固定 VAE、两个文本编码器和主 UNet；阶段 2 和阶段 3 固定 $\Phi_R$，阶段 2 额外固定 $A_\psi$。}
\label{tab:training_hyperparams}
{\scriptsize
\setlength{\tabcolsep}{2.5pt}
\resizebox{\textwidth}{!}{%
\begin{tabular}{p{0.055\textwidth}p{0.15\textwidth}p{0.14\textwidth}p{0.10\textwidth}p{0.07\textwidth}p{0.11\textwidth}p{0.19\textwidth}p{0.08\textwidth}p{0.085\textwidth}}
\toprule
阶段 & 初始化来源 & 可训练模块 & 有效损失 & 学习率 & T/R/B/M & GPU$\times$每卡 batch$\times$累积 & 全局 batch & 优化步数 \\
\midrule
1 & RealVisXL V3.0、PhotoMaker v1 & $\Phi_R$、ControlNet、$A_\psi$ & 完整目标 & $5\times10^{-5}$ & .20/.20/.20/.40 & $6\times12\times1$（前 55,000 步）；$4\times12\times1$（后 90,900 步） & 72；48 & 145,900 \\
2 & 阶段 1 最终检查点 & ControlNet & $\mathcal{L}_{\mathrm{diff}}$ & $5\times10^{-5}$ & .15/.20/.15/.50 & $3\times16\times1$ & 48 & 43,770 \\
3 & 阶段 1 的 $\Phi_R$；阶段 2 的 ControlNet 与 $A_\psi$ & ControlNet、$A_\psi$ & 完整目标 & $1\times10^{-5}$ & .20/.20/.15/.45 & $3\times12\times1$ & 36 & 19,450 \\
\bottomrule
\end{tabular}%
}
}
\end{table*}

阶段 1 共处理 8,323,200 个样本实例：前 55,000 步使用全局 batch 72，随后 90,900 步使用全局 batch 48，按 70,032 条训练记录计算约为 118.85 次数据遍历。阶段 2 和阶段 3 分别执行 43,770 和 19,450 个优化器更新步。除 ConsistentFace-Global、阶段比较和训练策略消融外，正文中的完整模型结果均取自阶段 3 最终检查点。

三个阶段均使用 bf16 精度和 AdamW 优化器，$\beta_1=0.9$、$\beta_2=0.999$、$\epsilon=10^{-8}$、权重衰减为 $10^{-2}$。学习率采用带 500 个预热步的常数调度，梯度累积步数为 1，不使用梯度裁剪。其余公共参数见表~\ref{tab:common_hyperparams}。推理固定噪声时间步为 799，采用 40 步去噪、CFG 4.0、ControlNet 条件强度 1.1 和小波颜色校正。

\begin{table*}[!t]
\centering
\caption{模型与训练的公共超参数。}
\label{tab:common_hyperparams}
{\scriptsize
\setlength{\tabcolsep}{4pt}
\resizebox{\textwidth}{!}{%
\begin{tabular}{p{0.19\textwidth}p{0.25\textwidth}p{0.19\textwidth}p{0.25\textwidth}}
\toprule
参数 & 设置 & 参数 & 设置 \\
\midrule
输入分辨率 & $512\times512$ & 主训练随机种子 & 233 \\
条件 token 上限 & 8 & 条件嵌入维度 & 2048 \\
适配器隐藏维度 & 512 & 门控隐藏维度 & 128 \\
排序间隔 & $m_g=m_a=0.20$ & 损失权重 & $\lambda_g=\lambda_a=\lambda_b=0.05$，$\lambda_d=0.02$ \\
数据加载失败重试 & 50 & 梯度检查点 & 关闭 \\
数值精度 & bf16 & 身份条件屏蔽概率 & 阶段 1：0；阶段 2/3：0.10 \\
\bottomrule
\end{tabular}%
}
}
\end{table*}

为使条件分配方法可直接比较，同一主干实验固定 SDXL/ControlNet 恢复主干 \cite{podell2023sdxl,zhang2023controlnet}、参考编码器、训练清单和推理参数。所有可训练双先验变体均接收 $(y,\mathcal{R},T)$，并使用五个随机种子训练和评测。

\subsubsection{测试数据与冲突构造}
测试集同时覆盖合成退化和真实退化。合成部分沿用 FaceMe 的 LFW-Test、WebPhoto-Test 和 WIDER-Test 划分 \cite{liu2025faceme}，并分别构造 1,711/1,707、407/407 和 970/969 条匹配/冲突记录；所有记录均按照可用身份标识与训练集去重。真实退化部分包含 640 张低照度或社交媒体图像（320 个身份）和 520 张老照片或强压缩图像（260 个身份）。每个身份均配有同身份参考图像，文本由注释者在不查看恢复结果的条件下撰写，标注分歧通过第二轮复核解决。

合成退化包括模糊、下采样、噪声和 JPEG 压缩 \cite{wang2021gfpgan,zhou2022codeformer,wang2023dr2}。在保持目标身份和可靠退化观测不变的条件下，测试记录改变参考集合或文本提示，形成匹配条件、轻度与重度文本冲突、文本错误、过时参考、错误参考和双条件失效。错误文本取自其他身份，不参与阈值或检查点选择。自由文本子集进一步覆盖跨年龄参考、遮挡、少见配饰和多个局部属性同时冲突。

错误来源实验额外构造配对反事实测试集，使文本错误与参考错误样本的 $s_{RT}$ 和退化强度 $q$ 分布相匹配，并平衡退化类型、参考数量、参考质量、身份划分和提示长度。由此可在跨先验相似度接近时单独检验观测核验和参考集稳定性的作用。所有冲突标签仅用于训练监督和评测分组，推理阶段均不提供。

\subsubsection{比较方法与评价指标}
实验包含两类比较。外部比较覆盖通用盲人脸恢复、参考驱动恢复和文本驱动恢复，用于展示不同输入接口下的性能范围；同一主干比较固定输入、参数规模和训练预算，只改变条件融合或可信度分配模块，用于检验本文机制本身。

恢复质量使用 PSNR、SSIM \cite{wang2004ssim} 和 LPIPS \cite{zhang2018lpips}。身份保持由 ArcFace、AdaFace、MagFace 和 CurricularFace 评测 \cite{deng2019arcface,kim2022adaface,meng2021magface,huang2020curricularface}，主表报告四个识别器的平均身份差距；该指标衡量同一样本在匹配条件与冲突条件下的身份相似度下降。外观一致性使用 CelebA 属性分类器 \cite{liu2015faceattributes}、OpenCLIP 属性评分 \cite{radford2021clip,ilharco2021openclip} 和开放属性解析器。

为补充平均相似度，本文还报告阈值级验证和身份检索。前者判断恢复结果与同身份参考的相似度是否超过固定阈值，并以 TAR 降幅描述冲突造成的验证性能下降；后者将恢复结果作为查询，统计身份库中的 Rank-1 和 Rank-5 命中率。排名反转率则统计传统画质指标选出的结果在人脸身份或目标属性上反而更差的比例。公式和判定细节见附录~\ref{app:metric_definitions}；所有阈值均由验证集确定，并在测试阶段固定。

\subsection{与现有方法比较}
\subsubsection{不同先验输入设置}
表~\ref{tab:external_context} 首先比较采用不同先验输入的代表性方法。该表反映从仅使用退化观测，到加入文本、参考以及双先验后的性能变化；由于输入接口和恢复主干不同，数值用于说明任务范围，不作为可信度分配机制的直接对照。

\begin{table*}[!t]
\centering
\caption{不同输入条件下的结果概览。}
\label{tab:external_context}
{\scriptsize
\renewcommand{\arraystretch}{0.92}
\setlength{\tabcolsep}{1mm}
\resizebox{\textwidth}{!}{%
\begin{tabular}{p{0.28\textwidth}cccccc}
\toprule
方法 & 先验 & PSNR $\uparrow$ & LPIPS $\downarrow$ & ArcFace $\uparrow$ & Major Attr-F1 $\uparrow$ & 文本错误身份差距 $\downarrow$ \\
\midrule
GFPGAN~\cite{wang2021gfpgan} & $y$ & 24.62 & 0.524 & 0.462 & -- & -- \\
CodeFormer~\cite{zhou2022codeformer} & $y$ & 25.14 & 0.506 & 0.481 & -- & -- \\
RestoreFormer++~\cite{wang2023restoreformerpp} & $y$ & 25.31 & 0.497 & 0.493 & -- & -- \\
DiffBFR~\cite{qiu2023diffbfr} & $y$ & 25.46 & 0.486 & 0.501 & -- & -- \\
OSDFace~\cite{wang2025osdface} & $y$ & 25.72 & 0.472 & 0.514 & -- & -- \\
DynFaceRestore~\cite{do2025dynfacerestore} & $y$ & 25.80 & 0.468 & 0.518 & -- & -- \\
FLIPNET~\cite{miao2025flipnet} & $y$ & 25.84 & 0.464 & 0.520 & -- & -- \\
TBFR~\cite{an2026tbfr} & $y,T$ & 25.58 & 0.476 & 0.498 & 0.706 & 0.042 \\
MCS~\cite{li2025mcs} & $y,T$ & 25.74 & 0.469 & 0.506 & 0.724 & 0.038 \\
A$^2$BFR~\cite{zhu2026a2bfr} & $y,T$ & 25.77 & 0.466 & 0.510 & 0.742 & 0.036 \\
FaceMe~\cite{liu2025faceme} & $y,\mathcal{R}$ & 25.88 & 0.461 & 0.535 & 0.514 & 0.022 \\
ConsistentFace-Global & $y,\mathcal{R},T$ & 26.04 & 0.452 & 0.548 & 0.758 & 0.014 \\
\bottomrule
\end{tabular}%
}
}
\end{table*}

表中 ConsistentFace-Global 为仅含全局分配的早期版本，报告外部固定清单上的单种子结果。它说明双先验输入在该设置下具有可用性；完整模型的效果及各组件贡献由下述同一主干五种子实验检验。由于模型版本和评测切片不同，两表的绝对数值不直接比较。

\subsubsection{同一主干比较}
\paragraph{基线构造。}
基线按条件组合方式分为三组。直接融合组包括 Token 级融合和参考--文本分支适配，前者直接组合全部条件 token，后者在个性化参考通路上增加可训练语言分支。单一或组合证据组包括图文嵌入一致性分配器、CLIP/ArcFace 一致性分配器和不确定性校准分配器，分别依据图文相似度、语义与身份一致性以及条件不确定性调整权重。结构消融组包括去除排序损失的多源证据分配器、无回退分配器和 RT-only；三者分别移除排序监督、回退 token 或除 $s_{RT}$ 外的证据。

所有基线采用相同输入、优化流程和推理设置，可训练参数规模与 ConsistentFace 保持一致；扩散主干、VAE、文本编码器、参考编码器和采样器均固定。检查点由预先划分的验证集确定，测试切片不参与模型选择。

\paragraph{五个随机种子下的结果。}
\begin{table*}[!t]
\centering
\caption{同一主干下的五种子平均结果。$\Delta$ 表示相对完整模型的差距。}
\label{tab:main_multiseed}
{\scriptsize
\renewcommand{\arraystretch}{0.92}
\setlength{\tabcolsep}{0.7mm}
\resizebox{\textwidth}{!}{%
\begin{tabular}{p{0.29\textwidth}cccccc}
\toprule
方法 & 平均冲突差距 $\downarrow$ & Major Attr-F1 $\uparrow$ & 平均 TAR 降幅 $\downarrow$ & 反转 $\downarrow$ & $\Delta$Gap & $\Delta$反转（百分点） \\
\midrule
Token 级交叉注意力融合 & 0.0154 & 0.701 & 0.034 & 55.8\% & 0.0045 & 11.2 \\
图文嵌入一致性分配器 & 0.0156 & 0.698 & 0.035 & 56.4\% & 0.0047 & 11.8 \\
参考--文本分支适配 & 0.0158 & 0.707 & 0.035 & 55.1\% & 0.0049 & 10.5 \\
CLIP/ArcFace 一致性分配器 & 0.0150 & 0.709 & 0.032 & 54.2\% & 0.0041 & 9.6 \\
不确定性校准分配器 & 0.0151 & 0.704 & 0.033 & 54.7\% & 0.0042 & 10.1 \\
去除排序损失的多源证据分配器 & 0.0143 & 0.714 & 0.031 & 52.6\% & 0.0034 & 8.0 \\
无回退分配器 & 0.0147 & 0.706 & 0.032 & 53.9\% & 0.0038 & 9.3 \\
ConsistentFace RT-only & 0.0118 & 0.724 & 0.026 & 47.0\% & 0.0009 & 2.4 \\
ConsistentFace 完整版 & \textbf{0.0109} & \textbf{0.729} & \textbf{0.023} & \textbf{44.6\%} & -- & -- \\
\bottomrule
\end{tabular}%
}
}
\end{table*}

表~\ref{tab:main_multiseed} 中，完整模型在四项指标上均优于直接融合、证据分配基线和结构消融变体。与无回退分配器相比，身份差距和 TAR 降幅分别降低 $0.0038$ 和 $0.009$，排名反转率降低 9.3 个百分点，Major Attr-F1 提高 $0.023$，说明外部条件均不可信时保留独立恢复路径是必要的。相较 RT-only，完整证据进一步改善身份差距、Major Attr-F1、TAR 降幅和排名反转率，表明仅依赖参考--文本相似度无法充分支持条件选择。去除排序损失后四项指标均退化，则说明显式分配监督能够稳定证据到权重的映射。

\subsubsection{强双先验基线}
同一主干比较之外，表~\ref{tab:strong_dual_prior} 进一步采用 IP-Adapter~\cite{ye2023ipadapter}、PhotoMaker~\cite{li2024photomaker}、InstantID~\cite{wang2024instantid} 和 PuLID~\cite{guo2024pulid} 构建双先验基线，用于检验结果是否仅来自更强的图像提示或身份定制模块。

\begin{table*}[!t]
\centering
\caption{强双先验基线的定量比较。ConsistentFace 行复用表~\ref{tab:main_multiseed} 的五种子均值。}
\label{tab:strong_dual_prior}
{\scriptsize
\renewcommand{\arraystretch}{0.92}
\setlength{\tabcolsep}{1mm}
\resizebox{\textwidth}{!}{%
\begin{tabular}{p{0.40\textwidth}cccc}
\toprule
双先验基线 & 差距 $\downarrow$ & Major Attr-F1 $\uparrow$ & 错误属性采纳率 $\downarrow$ & 反转 $\downarrow$ \\
\midrule
ControlNet + IP-Adapter + 文本分支 & 0.0155 & 0.719 & 0.240 & 55.4\% \\
PhotoMaker + 恢复适配器 & 0.0149 & 0.707 & 0.246 & 54.8\% \\
InstantID + ControlNet 恢复代理 & 0.0136 & 0.699 & 0.268 & 52.7\% \\
PuLID + ControlNet 恢复代理 & 0.0132 & 0.704 & 0.259 & 52.1\% \\
PhotoMaker + 文本分支 & 0.0146 & 0.723 & 0.226 & 53.5\% \\
ConsistentFace 完整版 & \textbf{0.0109} & \textbf{0.729} & \textbf{0.158} & \textbf{44.6\%} \\
\bottomrule
\end{tabular}%
}
}
\end{table*}

ConsistentFace 在四项指标上均取得最佳平均结果。InstantID 和 PuLID 式基线的身份差距分别为 $0.0136$ 和 $0.0132$，ConsistentFace 将其降至 $0.0109$，同时将错误属性采纳率降至 $0.158$。这组结果表明，增强身份条件本身不能替代冲突条件下的可信度判断。

\subsection{消融实验}
\subsubsection{局部属性分配与回退}
表~\ref{tab:slot_ablation} 从仅使用全局分配开始，依次加入回退条件、局部属性分配和完整多源证据，用于分析整体条件选择与局部冲突处理之间的关系。局部 Target Attr-F1 和错误属性采纳率分别衡量目标属性保留与错误外观进入恢复结果的程度。

\begin{table*}[!t]
\centering
\caption{局部属性分配与回退条件消融。}
\label{tab:slot_ablation}
{\scriptsize
\renewcommand{\arraystretch}{0.92}
\setlength{\tabcolsep}{1mm}
\resizebox{\textwidth}{!}{%
\begin{tabular}{p{0.32\textwidth}ccccc}
\toprule
变体 & 差距 $\downarrow$ & Major Attr-F1 $\uparrow$ & 局部 Target Attr-F1 $\uparrow$ & 错误属性采纳率 $\downarrow$ & 反转 $\downarrow$ \\
\midrule
仅全局分配器 & 0.0134 & 0.711 & 0.728 & 0.231 & 50.8\% \\
全局分配器 + 回退 & 0.0127 & 0.716 & 0.739 & 0.211 & 49.2\% \\
RT-only + 局部属性分配 & 0.0118 & 0.724 & 0.756 & 0.178 & 47.0\% \\
多源证据 + 局部属性分配 & \textbf{0.0109} & \textbf{0.729} & \textbf{0.762} & \textbf{0.158} & \textbf{44.6\%} \\
\bottomrule
\end{tabular}%
}
}
\end{table*}

仅使用全局分配时，局部 Target Attr-F1 为 $0.728$，错误属性采纳率为 $0.231$；加入回退条件后两项指标分别改善至 $0.739$ 和 $0.211$。在此基础上，RT-only 局部属性分配将局部 Target Attr-F1 提高到 $0.756$，并将错误属性采纳率降至 $0.178$，说明属性级分配能够减少局部错误扩散。完整多源证据进一步把两项指标改善至 $0.762$ 和 $0.158$。最后一步同时改变证据配置和局部分配输入，因此该结果支持二者的联合设计，不将全部增益单独归因于属性分组。

\subsubsection{参考与文本错误来源的区分}
本节检验多源证据能否在跨先验相似度接近时区分错误来源。配对测试匹配 $s_{RT}$ 和 $q$ 的分布，并分别以文本错误和参考错误为正类构造两个 one-vs-rest 判别任务；由此，完整证据相对 $s_{RT}$ 与 $q$ 的增益主要来自观测核验分数 $s_{yR}$、$s_{yT}$ 和参考集稳定性 $s_{RR}$。

图~\ref{fig:counterfactual_errors} 给出两个配对案例。两者具有接近的 $s_{RT}$ 与 $q$，但文本错误样本中的观测与参考更加一致，分配器因而提高参考偏好；参考错误样本中的观测与文本更加一致，分配器转而提高文本贡献。

\begin{figure*}[!t]
\centering
\includegraphics[width=\textwidth]{figures/fig05_counterfactual_errors.png}
\caption{配对反事实条件下的错误来源区分。两类样本具有匹配的参考--文本相似度和退化强度，但在外部条件与可靠观测的一致性、参考集内部稳定性上存在差异，因此完整证据能够区分文本错误与参考错误，并输出方向相反的条件偏好。}
\label{fig:counterfactual_errors}
\end{figure*}

\begin{table}[!t]
\centering
\caption{配对反事实错误来源判别的 AUROC。}
\label{tab:source_disambiguation}
\small
\setlength{\tabcolsep}{4pt}
\begin{tabular}{lccc}
\toprule
证据输入 & 文本错误 & 错误参考 & 平均 \\
\midrule
$s_{RT}$ 与 $q$ & 0.55 & 0.53 & 0.54 \\
完整证据向量 & \textbf{0.82} & \textbf{0.84} & \textbf{0.83} \\
\bottomrule
\end{tabular}
\end{table}

表~\ref{tab:source_disambiguation} 表明，仅使用参考--文本相似度和退化强度时，两类错误的平均 AUROC 为 0.54，接近随机判别；加入观测核验和参考集稳定性后，文本错误与参考错误的 AUROC 分别提高到 0.82 和 0.84。该结果说明，多源证据补充了 $s_{RT}$ 无法提供的错误方向信息。

\subsection{鲁棒性与进一步分析}
\subsubsection{跨评估器与人工评审}
属性结论使用 CelebA 分类器、OpenCLIP、开放属性解析器和独立人工子集交叉验证。表~\ref{tab:evaluator_robustness} 报告各评估器的属性分数，以及相对主评估器方法排序的 Spearman 相关系数。

\begin{table}[!t]
\centering
\caption{跨属性评估器结果。$\rho$ 表示与主评估器排序的 Spearman 相关系数。}
\label{tab:evaluator_robustness}
{\footnotesize
\setlength{\tabcolsep}{1mm}
\begin{tabular}{p{0.42\linewidth}ccc}
\toprule
评估器 & Major Attr-F1 & 局部 Attr-F1 & $\rho$ \\
\midrule
CelebA 属性分类器 & 0.729 & 0.762 & 1.00 \\
OpenCLIP 属性评分 & 0.718 & 0.748 & 0.87 \\
开放属性解析器 & 0.724 & 0.755 & 0.90 \\
人工评审子集 & 0.742 & 0.771 & 0.92 \\
\bottomrule
\end{tabular}
}
\end{table}

三种替代评估方式与主评估器的排序相关系数为 $0.87$ 至 $0.92$，说明属性结论不依赖单一分类器。人工身份评审另覆盖 900 个输出对，并在合成退化与真实退化切片间保持平衡。合成切片以配准清晰目标为身份锚点，真实切片以独立同身份参考为锚点；每对由三名注释者在隐藏提示和方法名称、随机左右顺序的界面上比较匹配提示与文本错误结果。“不确定”占原始投票的 8.7\%，未形成非不确定多数的样本仅用于一致性统计，偏好率以其余样本为分母。匹配提示结果的身份偏好率为 71.4\%，与 ArcFace 判断的一致率为 88.9\%，Fleiss $\kappa$ 为 0.62。表~\ref{tab:evaluator_robustness} 中的人工属性子集与该身份偏好任务相互独立。

\subsubsection{自然语言和自由文本提示}
表~\ref{tab:natural_breakdown} 将评测扩展到预设属性组外属性、跨年龄参考、复杂遮挡、混合冲突和长自由文本，比较无回退分配器与完整模型。该实验检验固定属性分组和模板化训练之外的提示变化。

\begin{table}[!t]
\centering
\caption{自然语言和自由文本提示切片。N/F 分别表示 NoFB 与完整模型，$n$ 为提示--输出对数。}
\label{tab:natural_breakdown}
{\scriptsize
\setlength{\tabcolsep}{1pt}
\begin{tabular}{p{0.24\linewidth}cccc}
\toprule
切片 & $n$ & 差距 N/F $\downarrow$ & Target F1 $\uparrow$ & 反转 N/F $\downarrow$ \\
\midrule
预设属性组外属性 & 260 & 0.0205/\textbf{0.0169} & 0.736 & 60.4\%/\textbf{54.7\%} \\
跨年龄参考 & 300 & 0.0178/\textbf{0.0129} & 0.771 & 57.2\%/\textbf{47.2\%} \\
面具或手部遮挡 & 240 & 0.0201/\textbf{0.0158} & 0.752 & 59.5\%/\textbf{51.9\%} \\
混合冲突 & 420 & 0.0193/\textbf{0.0147} & 0.760 & 58.8\%/\textbf{50.3\%} \\
长自由文本提示 & 500 & 0.0173/\textbf{0.0131} & 0.776 & 55.9\%/\textbf{46.1\%} \\
\bottomrule
\end{tabular}
}
\end{table}

完整模型在五个切片上均取得更低的身份差距和排名反转率。改进在跨年龄参考和长自由文本中最为明显，排名反转率分别降低 10.0 和 9.8 个百分点；在预设属性组外属性上，完整模型仍将身份差距从 $0.0205$ 降至 $0.0169$。这些结果说明回退条件不仅作用于预设属性组内的冲突，也能缓解无法由局部分组完整描述的条件失效。

\subsubsection{真实退化与退化变化压力测试}
在 640 张低照度或社交媒体图像和 520 张老照片或强压缩图像上，错误文本相对匹配文本的 ArcFace 相似度分别下降 $0.0318$ 和 $0.0285$。真实退化切片没有配准清晰目标，因此本节不以 PSNR、SSIM 或 LPIPS 推断身份保持程度，而采用同身份参考下的识别器相似度和人工成对判断。

为检验上述现象是否局限于训练退化，压力测试进一步覆盖运动模糊、局部遮挡、低照度、强 JPEG 压缩和多平台压缩。测试文本由人工撰写，退化类型不作为分配器的训练标签。

\begin{table*}[!t]
\centering
\caption{退化变化压力测试。$\Delta$ ArcFace 表示“错误文本减匹配文本”，配对 $p$ 值经 Holm 校正。}
\label{tab:ood_degradation}
\small
\setlength{\tabcolsep}{4pt}
\begin{tabular}{lcccc}
\toprule
退化类型 & 图像数 & $\Delta$ ArcFace & 95\% CI & 配对 $p$ \\
\midrule
运动模糊 & 420 & $-0.0294$ & $[-0.0336,-0.0251]$ & $<0.001$ \\
局部遮挡 & 360 & $-0.0267$ & $[-0.0312,-0.0225]$ & $<0.001$ \\
低照度 & 410 & $-0.0321$ & $[-0.0369,-0.0278]$ & $<0.001$ \\
强 JPEG 压缩 & 450 & $-0.0279$ & $[-0.0320,-0.0235]$ & $<0.001$ \\
多平台压缩 & 380 & $-0.0306$ & $[-0.0352,-0.0260]$ & $<0.001$ \\
\bottomrule
\end{tabular}
\end{table*}

五类退化上的置信区间均未跨越零，说明错误文本造成的身份相似度下降在不同退化下持续存在。低照度条件下降幅最大，多平台压缩次之，局部遮挡的降幅最小。该结果将正文结论从训练所用退化扩展到未作为路由标签的退化变化。

\section{局限性与伦理说明}
\label{sec:main_limitations}

ConsistentFace 的条件判断依赖退化观测中仍可辨认的信息。当退化极强或关键区域被遮挡时，观测可能不足以区分两类外部条件，分配结果因而仍会出错。当前五类属性组也不能完整覆盖表情、妆容、细粒度肤色、脸型、痣疤、复杂光照或多人交互等开放属性。完整模型的排名反转率仍为 44.6\%，说明传统画质指标与身份、属性评价之间仍存在明显差异。后续可进一步研究开放属性解析、遮挡感知身份置信度、时间感知参考加权和不确定性校准。

此外，可信度分配不等同于身份认证。生成式人脸恢复依赖生成先验补充退化输入中难以辨认的细节 \cite{wang2021gfpgan,zhou2022codeformer,lin2024diffbir}，输出不应被视为原始细节的事实性还原或身份核验证据。ConsistentFace 仅适用于经授权的照片修复、媒体增强和恢复算法评测，不建议用于门禁、执法、金融开户、招聘筛查等高风险场景。面部识别系统还伴随隐私和准确性风险 \cite{nist2019frt,gao2020facialrecognition}。

模型或数据发布还应附带输入退化记录、参考来源摘要、可信度诊断和失败模式说明。本文的真实退化切片仅用于压力诊断；发布相应结果时，应同时给出数据来源、授权范围、身份级划分清单和去重哈希。对于受许可、肖像权或隐私约束的人脸数据，宜只发布提示、评测清单、哈希值、聚合统计和可复现实验脚本；涉及未成年人、敏感身份或不可撤回同意的数据不应随模型发布原始图像。

\section{结论}
\label{sec:main_conclusion}

本文研究参考图像与文本提示不一致时的身份保持人脸恢复，并提出可信度分配适配器 ConsistentFace。该方法以可靠退化观测为锚点，从观测、参考集和文本之间的一致性及退化强度中构造分配证据，再通过全局与局部两级机制选择参考、文本和回退条件；推理过程无须冲突标签。同一主干比较说明可信度分配优于直接融合，消融与反事实实验进一步验证了多源证据、局部属性分配和回退条件的作用。自然语言、多评估器、人工评审和退化变化测试表明，这些改进在不同提示、评估方式和退化类型下保持一致。

\appendices
\section{指标定义与统计协议}
\label{app:metric_definitions}

给定身份评估器 $h$ 的相似度 $S_h$，身份差距定义为
\begin{equation}
\begin{aligned}
\mathrm{IdentityGap}=\mathbb{E}_{h}\big[
&\mathbb{E}_{\mathrm{match}}S_h(\hat{x}^{\mathrm{match}},x)\\
&-\mathbb{E}_{\mathrm{conf}}S_h(\hat{x}^{\mathrm{conf}},x)
\big].
\end{aligned}
\end{equation}
其中，$\mathrm{match}$ 和 $\mathrm{conf}$ 分别表示匹配条件和对应冲突条件。对于候选集 $\mathcal{C}_i$，排名反转定义为
\begin{equation}
\begin{aligned}
\mathrm{Rev}_i
&=\mathbf{1}\{\exists j\in\mathcal{C}_i:\ H(k^\star)<H(j)-\tau\},\\
k^\star
&=\arg\max_{k\in\mathcal{C}_i}M(k),
\end{aligned}
\end{equation}
其中，$M\in\{\mathrm{PSNR},\mathrm{SSIM},-\mathrm{LPIPS}\}$ 为画质排序指标，$H$ 为身份或目标属性指标，$\tau=10^{-4}$ 为验证集上固定的平局边距。若画质指标选出的 $k^\star$ 在 $H$ 上低于任一候选超过 $\tau$，则记为一次排名反转。属性组 $a$ 上的错误属性采纳量定义为
\begin{equation}
\mathrm{WPA}_a=\mathbf{1}\{E_a(\hat{x})=z_a^{\mathrm{wrong}}\land E_a(\hat{x})\ne z_a^{\mathrm{target}}\}.
\end{equation}
其中，$E_a$ 为属性 $a$ 的评估器，$z_a^{\mathrm{wrong}}$ 和 $z_a^{\mathrm{target}}$ 分别为冲突条件给出的错误属性和目标属性。错误属性采纳率为 $\mathrm{WPA}_a$ 在样本与属性组上的平均值。

阈值级验证沿用人脸验证中的固定阈值判定。给定验证集上由目标误接受率 $\eta$ 确定的阈值 $\delta_{\eta}$，若恢复结果与同身份参考的相似度不低于该阈值，则判定为通过验证。真接受率（TAR）定义为
\begin{equation}
\mathrm{TAR}_{\eta}
=
\frac{1}{|\mathcal{P}^{+}|}
\sum_{(i,j)\in\mathcal{P}^{+}}
\mathbf{1}\{S_h(\hat{x}_i,r_j)\geq \delta_{\eta}\},
\end{equation}
其中，$\mathcal{P}^{+}$ 为同身份验证对，$r_j$ 为同身份参考图像。本文报告 TAR@FAR=$10^{-3}$，并用匹配提示与冲突提示之间的差值表示 TAR 降幅。等错误率（EER）为误接受率与误拒绝率相等时的错误率；数值越低表示验证边界越稳定。

检索指标把恢复结果作为查询，将候选库中同身份图像视为正确项。设 $\operatorname{rank}_i$ 为正确身份在相似度排序中的最高名次，则 Rank-$k$ 检索准确率定义为
\begin{equation}
\mathrm{Rank}\text{-}k
=
\frac{1}{N}
\sum_{i=1}^{N}
\mathbf{1}\{\operatorname{rank}_i\leq k\}.
\end{equation}
Rank-1 衡量首位检索是否命中正确身份，Rank-5 衡量前五个候选中是否包含正确身份。

以上阈值均由验证集确定，并在测试阶段保持不变。五种子实验报告均值；置信区间通过身份级配对 bootstrap 计算，多重比较采用 Holm 校正。

\section{补充实验}
\subsection{简化全局路由对照}
为分离全局路由与局部属性分配的作用，本文在较小的同一主干切片上使用仅保留整体一致性证据 $s$ 和退化强度 $q$ 的简化版本。由于该切片和模型配置不同于正文主表，表~\ref{tab:global_route_controls} 只用于比较路由策略，不直接比较绝对数值。

\begin{table}[!h]
\centering
\caption{简化全局路由切片对照。}
\label{tab:global_route_controls}
{\footnotesize
\setlength{\tabcolsep}{2pt}
\begin{tabular}{lccc}
\toprule
方法 & 路由监督 & 身份差距 $\downarrow$ & Major F1 $\uparrow$ \\
\midrule
冻结 FaceMe + 文本注入 & 无 & 0.0258 & 0.604 \\
固定拼接 & 无 & 0.0194 & 0.662 \\
固定相加 & 无 & 0.0190 & 0.668 \\
可学习非路由门控 & 无 & 0.0176 & 0.684 \\
仅一致性门控 & 无 & 0.0171 & 0.690 \\
随机路由标签 & 随机 & 0.0231 & 0.641 \\
错误路由标签 & 错误 & 0.0273 & 0.618 \\
ConsistentFace-Global & 排序 & 0.0137 & 0.710 \\
Oracle 路由 & 真实路由 & \textbf{0.0115} & \textbf{0.721} \\
\bottomrule
\end{tabular}
}
\end{table}

固定融合与可学习门控均能减小身份差距，但随机或错误路由标签会明显削弱性能。ConsistentFace-Global 的身份差距低于固定融合和无监督门控，同时与 Oracle 路由仍有 $0.0022$ 的差距，说明学习得到的全局分配有效，但尚未达到已知真实路由时的上限。

\subsection{多识别器与验证指标}
表~\ref{tab:identity_verification} 从相似度、人工偏好、阈值验证和身份检索四个角度比较匹配文本与错误文本。四个身份识别器分别报告余弦相似度，TAR 与 EER 描述固定验证阈值下的判定变化，Rank-1 和 Rank-5 描述身份库检索结果。

\begin{table}[!t]
\centering
\caption{匹配文本与错误文本下的身份验证。}
\label{tab:identity_verification}
{\footnotesize
\setlength{\tabcolsep}{3pt}
\begin{tabular}{lccc}
\toprule
指标 & 匹配文本 & 错误文本 & 变化 \\
\midrule
ArcFace 相似度 $\uparrow$ & 0.5364 & 0.5231 & $-0.0133$ \\
AdaFace 相似度 $\uparrow$ & 0.5480 & 0.5337 & $-0.0143$ \\
MagFace 相似度 $\uparrow$ & 0.5315 & 0.5188 & $-0.0127$ \\
CurricularFace 相似度 $\uparrow$ & 0.5442 & 0.5310 & $-0.0132$ \\
人工成对身份偏好 $\uparrow$ & 71.4\% & 28.6\% & $-42.8$ pp \\
TAR@FAR=$10^{-3}$ $\uparrow$ & 0.742 & 0.704 & $-0.038$ \\
EER $\downarrow$ & 0.086 & 0.098 & $+0.012$ \\
Rank-1 检索 $\uparrow$ & 78.3\% & 74.2\% & $-4.1$ pp \\
Rank-5 检索 $\uparrow$ & 91.4\% & 88.6\% & $-2.8$ pp \\
\bottomrule
\end{tabular}
}
\end{table}

四个识别器在错误文本下均给出 $0.0127$ 至 $0.0143$ 的相似度下降；TAR、EER 和检索准确率也沿同一方向变化。该结果说明身份下降并非由单一 ArcFace 评估器造成，人工判断、阈值验证与检索指标为平均相似度提供了互补证据。
\bibliographystyle{IEEEtran}
\bibliography{references}

\end{document}
